% ------------------------------------------------------------------------------
% 2 Data  3 Methods
% ------------------------------------------------------------------------------
\section{Data and Endpoints}\label{sec:data}

\subsection{Source and patients}

The analysis uses the individual-patient data of the NCCTG adjuvant colon
cancer trial as distributed with the R package \texttt{survival}
\citep{therneau2026}: 929 patients with resected stage B/C colon
carcinoma, all from a single study, randomized 1:1:1 to Observation
($n=315$), Levamisole ($n=310$), or Levamisole plus 5-FU ($n=304$). The
recorded baseline covariates are age, sex, bowel obstruction,
perforation, adherence to nearby organs, number of positive lymph nodes,
tumour differentiation, extent of local spread, and the
surgery-to-registration interval. Randomization produced well-balanced
arms: across eleven baseline variables the only nominal imbalance was
sex (53.6\% female in Lev+5FU vs.\ 47.3\% in Observation; three-arm
global exact test $p=0.027$, pairwise $p=0.13$), unremarkable among
eleven simultaneous tests and not surviving multiplicity correction;
pairwise standardized mean differences never exceeded $|0.216|$, and the
overall picture is consistent with adequate randomization concealment.

\subsection{Endpoint construction}

The raw dataset stores two records per patient: \texttt{etype~=~1} records
time to tumour recurrence with death treated as censoring, and
\texttt{etype~=~2} records overall survival. Because current adjuvant
colorectal conventions define DFS as recurrence \emph{or} death from any
cause, the data-curation module constructs the composite DFS endpoint
explicitly (506 events; 468 recurrences plus 38 deaths without recorded
recurrence), while overall survival (452 deaths) serves as the secondary
endpoint and the recurrence-only coding with death as a competing event
(38 competing events) is retained for sensitivity analyses. Covariate
values are verified to be identical across the two records of each
patient; event times satisfy the required orderings (every recurrence
time does not exceed the corresponding death time, and deaths without
recurrence share the death time). Median follow-up, estimated by reverse
Kaplan--Meier, is 6.4 years (maximum 9.1 years).

\subsection{Missing data}

Two covariates have missing values: number of positive lymph nodes
(18 patients, 1.9\%) and tumour differentiation (23 patients, 2.5\%);
95.6\% of patients are complete cases. Missing values are handled by
multiple imputation as described below, with complete-case analyses
reported for comparison.

\begin{table}[htbp]
\centering\small
\caption{Baseline characteristics by randomized arm (percentages of the
full arm; nodes and differentiation include a small missing category).}
\label{tab:table1}
\begin{tabular*}{\textwidth}{@{\extracolsep{\fill}} l c c c}
\toprule
Characteristic & Observation & Levamisole & Levamisole + 5-FU \\
\midrule
Randomized ($n$)              & 315 & 310 & 304 \\
Age, years, mean (SD)         & 59.5 (12.0) & 60.1 (11.6) & 59.7 (12.3) \\
Male sex, $n$ (\%)            & 166 (52.7) & 177 (57.1) & 141 (46.4) \\
Obstruction, $n$ (\%)         & 63 (20.0)  & 63 (20.3)  & 54 (17.8)  \\
Perforation, $n$ (\%)         & 9 (2.9)   & 10 (3.2)   & 8 (2.6)   \\
Adherence to organs, $n$ (\%) & 47 (14.9)  & 49 (15.8)  & 39 (12.8)  \\
Positive nodes, median [IQR]  & 2 [1--5]   & 2 [1--5]   & 2 [1--4]   \\
$\geq 4$ positive nodes, $n$ (\%) & 119 (37.8) & 105 (33.9) & 92 (30.3) \\
Poor differentiation, $n$ (\%)& 52 (16.5)  & 44 (14.2)  & 54 (17.8)  \\
Serosal/contiguous extent, $n$ (\%) & 269 (85.4) & 271 (87.4) & 262 (86.2) \\
Long surgery-to-registration, $n$ (\%) & 224 (71.1) & 230 (74.2) & 228 (75.0) \\
\midrule
Composite DFS events, $n$ (\%) & 190 (60.3) & 182 (58.7) & 134 (44.1) \\
Deaths, $n$ (\%)                & 168 (53.3) & 161 (51.9) & 123 (40.5) \\
\bottomrule
\end{tabular*}
\end{table}

\section{Methods}\label{sec:methods}

\subsection{Frequentist efficacy analysis}\label{sec:methods-freq}

Disease-free survival and overall survival were summarized by
Kaplan--Meier curves with pointwise 95\% confidence intervals (log--log
transformation) and compared with log-rank tests, both overall (2 degrees
of freedom) and pairwise between arms. Treatment effects were estimated
with Cox proportional-hazards models fitted by Efron's method for ties.
Three specifications are reported for the primary endpoint: an unadjusted
model; a covariate-adjusted complete-case model; and the pre-specified
covariate-adjusted model fitted on each of the $m=20$ multiply imputed
datasets and pooled by Rubin's rules \citep{rubin1987}. The imputation
models used chained equations with the Nelson--Aalen cumulative hazard
and the event indicator as predictors, following
\citet{white2009} and \citet{mice2011}, so that imputation is
compatible with the subsequent Cox analysis. The adjustment set
comprised treatment arm, age, sex, obstruction, perforation, adherence,
number of positive nodes, differentiation, extent of local spread, and
surgery-to-registration interval; an analogous pooled model was fitted
for overall survival.

The proportional-hazards assumption was examined with scaled
Schoenfeld residuals \citep{grambsch1994}, both graphically (loess
smooths against time) and with the global chi-square test. Two
pre-specified sensitivity analyses address potential violations and
endpoint-definition concerns: first, flexible functional forms for age
and nodes via restricted penalized splines (4 degrees of freedom each),
compared with the linear model by likelihood-ratio test; second,
competing-risk analyses on the recurrence endpoint, contrasting the
Aalen--Johansen cumulative incidence function
\citep{aalen1978} with the naive one-minus-Kaplan--Meier estimate, and
fitting Fine--Gray subdistribution-hazard models \citep{finegray1999}
adjusted for the same covariate set.

\subsection{Absolute benefit and restricted mean survival time}
\label{sec:methods-rmst}

Because hazard ratios are not directly interpretable as absolute
benefit, five-year Kaplan--Meier survival probabilities with 95\%
confidence intervals were computed by arm for both endpoints, along with
between-arm risk differences and numbers needed to treat (NNT) whose 95\%
confidence intervals were obtained by nonparametric bootstrap with
2{,}000 resamples.

In addition, the primary comparison was re-expressed through the
restricted mean survival time \citep{royston2011,uno2014}: the area
under the Kaplan--Meier survival curve up to a horizon $\tau$, estimated
at $\tau = 3,\ldots,8$ years with percentile bootstrap confidence
intervals (2{,}000 resamples), and summarized both as a between-arm
difference (disease-free years gained) and as a ratio. Unlike the hazard
ratio, the RMST difference is collapsible and requires no
proportional-hazards assumption. A covariate-adjusted RMST difference at
$\tau = 5$ years was obtained by regressing leave-one-out jackknife
pseudo-observations of $\rmst(5)$
\citep{andersen2003,parner2010} on treatment and the standardized
adjustment set, with bootstrap standard errors; because the estimand is
collapsible, the adjusted and unadjusted analyses target the same
quantity, providing a direct empirical contrast with the
non-collapsibility of the hazard ratio quantified in the simulation
study below. RMST differences were also computed within the obstruction
stratum, the covariate with the strongest proportional-hazards
violation.

\subsection{Subgroup and interaction analyses}

Seven clinically motivated subgroups were pre-specified for the primary
comparison of levamisole plus 5-FU against observation: age ($<60$
vs.\ $\geqslant 60$ years), sex, nodal burden ($<4$ vs.\ $\geq 4$
positive nodes), obstruction, local extent (submucosa/muscle vs.\
serosa/contiguous), differentiation (well/moderate vs.\ poor), and
surgery-to-registration interval. Within each subgroup level, unadjusted
Cox models estimated subgroup-specific hazard ratios; effect modification
was tested by Wald tests of treatment-by-subgroup interaction terms in
models adjusted for the nodal count. Given seven interactions tested,
only $p$-values below a Bonferroni-adjusted threshold of $0.05/7 \approx
0.0071$ are interpreted as evidence of modification.

\subsection{Trial redesign and re-enacted monitoring}
\label{sec:methods-design}

To redesign the trial under modern standards, planning assumptions were
taken from the re-analysis itself. The control arm's five-year DFS of
42.4\% and median DFS of 2.96 years imply an exponential baseline hazard
$\lambda_0 = \log 2 / 2.96 = 0.234$ per year; the clinically meaningful
target was set at a hazard ratio of $0.65$, close to the effect actually
observed. A fixed one-sided superiority design at $\alpha = 0.025$ with
90\% power, 1:1 allocation, three years of uniform accrual and a minimum
of five years of additional follow-up, requires 226 DFS events and
approximately 322 patients (Schoenfeld's formula
\citep{schoenfeld1983}, verified with \texttt{gsDesign::nSurv}
\citep{gsdesign2026}).

The modern redesign is a three-look group-sequential trial with analyses
after 50\% and 75\% of the planned events and at completion. Efficacy
boundaries use O'Brien--Fleming alpha-spending
\citep{o1979,landemets1983}; futility boundaries use a non-binding
Hwang--Shih--DeCani beta-spending with shape parameter $-4$
\citep{hwang1990}, which yields aggressive early stopping for futility
under the null while allowing the sponsor to disregard the futility
boundary without inflating the type-I error. The resulting maximum
sample size is 235 events (approximately 336 patients), a 4.1\%
inflation over the fixed design.

The monitoring plan was then \emph{re-enacted on the real trial data}:
because the public dataset contains no calendar dates, an
information-time snapshot of the interim stage was reconstructed by
administratively censoring every patient's follow-up at the time (years
since randomization) at which the 118th DFS event of the actual trial
was observed, refitting the Cox model on the censored-at-interim
snapshot, and comparing the resulting $Z$-statistic with the design
boundaries. Conditional power for the final analysis was computed under
three assumptions for the remaining drift (null effect, design
alternative, and the observed interim trend), with the sign convention
that positive $Z$ favours the experimental arm; the hand-rolled
calculation was cross-checked with \texttt{gsDesign}'s conditional-power
machinery.

\subsection{Simulation studies}\label{sec:methods-sim}

Three Monte Carlo studies quantify the operating characteristics behind
the design and analysis choices. The data-generating mechanism was
calibrated to the trial itself: covariates were drawn by bootstrap
resampling of the complete-case patients, the prognostic index used
hazard-ratio coefficients from the adjusted Cox model, event times were
exponential with hazard $\lambda_0 \exp(\mathrm{lp}_i + \log(\hr)
\cdot \mathrm{treat}_i)$, accrual times were uniform on three years, and
administrative censoring occurred at eight years.

\textbf{Study A (covariate-adjustment efficiency)} simulated 1{,}000
trials per scenario for a two-arm comparison with $n=322$ at hazard
ratios of $1.0$ (type-I error) and $0.72$ (power), crossing three levels
of prognostic strength obtained by scaling the calibrated covariate
coefficients by 0, 0.5, and 1 (standard deviations of the prognostic
index: 0, 0.208, and 0.416 on the log-hazard scale). Each replicate
fitted both the unadjusted and the fully adjusted Cox model, recording
estimates, model-based standard errors, empirical standard deviations,
and rejection frequencies. The contrast between the unadjusted estimator
and the data-generating conditional hazard ratio quantifies the
non-collapsibility of the hazard ratio \citep{martinussen2013}.

\textbf{Study B (sequential-design operating characteristics)}
simulated the full monitoring procedure of
Section~\ref{sec:methods-design} (3{,}000 trials per scenario) with
event-triggered looks and calendar-time censoring, under the alternative
($\hr=0.65$) and under the null with the futility boundary adhered to
and ignored.

\textbf{Study C (interval estimation under non-proportional hazards)}
simulated 500 trials with a piecewise-constant log-hazard ratio (HR
$0.45$ in year one, HR $0.80$ thereafter; the exponential
memorylessness property makes the piecewise construction exact),
comparing the coverage and width of Wald and percentile-bootstrap
confidence intervals for the log hazard ratio (100 bootstrap resamples
per replicate) against the simulation estimand.

\subsection{Machine-learning prediction}\label{sec:methods-ml}

Five prognostic models for the composite DFS endpoint were compared with
repeated ($5\times5$-fold) stratified cross-validation: a clinical Cox
model with main effects (benchmark); an elastic-net penalized Cox model;
a random survival forest; gradient boosting with Cox partial-likelihood
loss; and XGBoost with the Cox objective. Treatment assignment was
deliberately excluded from all predictors, so the models quantify
baseline risk only---a task distinct from treatment-effect estimation.
Discrimination was assessed with Harrell's C-index and the time-dependent
IPCW AUC at five years, prediction error with the integrated Brier score
over years 1--5 \citep{graf1999}, and calibration by deciles of
predicted risk against Kaplan--Meier observed risk. The apparent
(in-sample) performance of each model was also computed to quantify the
optimism gap. Permutation importance (20 shuffles per covariate)
interpreted the winning model.

\subsection{Bayesian re-analysis}\label{sec:methods-bayes}

A Weibull proportional-hazards model was fitted by NUTS Hamiltonian
Monte Carlo (PyMC; \citealp{salvatier2016}) to the primary two-arm
comparison, with right censoring handled by likelihood factorization
(the censored observations contribute their survival function). The
parametrization $H(t \mid x) = (t/\beta_0)^\alpha \exp(\eta_i)$, with
$\eta_i$ the linear predictor, makes $\exp(\beta_{\mathrm{treat}})$ the
hazard ratio. Three priors for the treatment log-hazard ratio were
compared: \emph{skeptical} ($\mathcal{N}(0, 0.421^2)$, placing 5\% prior
mass below $\hr = 0.5$); \emph{weakly-informative}
($\mathcal{N}(0, 1)$); and \emph{enthusiastic}
($\mathcal{N}(\log 0.65, 0.421^2)$, centred at the design alternative).
Sequential monitoring replayed evidence accumulation at 25\%, 50\%, 75\%,
and 100\% of the observed DFS events under the skeptical prior; posterior
predictive checks replicated the control arm's Kaplan--Meier curve from
100 posterior draws; and approximate leave-one-out cross-validation
\citep{vehtari2017} compared the Weibull baseline against the
exponential special case. Convergence was assessed by $\hat R$ and
effective sample sizes.

\subsection{Heterogeneous treatment effects}\label{sec:methods-cate}

Heterogeneity was characterized on the absolute five-year DFS event
risk. Jackknife pseudo-observations for the five-year risk
\citep{andersen2003,parner2010} converted the censored outcome into a
regression-ready continuous outcome. The population-average effect was
estimated three ways: the unadjusted Kaplan--Meier difference with
bootstrap confidence intervals; cross-fitted G-computation; and a
cross-fitted doubly robust AIPW estimator with influence-function
standard errors \citep{hernan2020}, using the empirical randomization
probability as the propensity. Conditional (patient-level) effects were
estimated by a cross-fitted T-learner (penalized outcome models by arm),
summarized by the distribution and quartiles of predicted absolute
benefit and by the dependence of benefit on baseline risk.

\subsection{Software and reproducibility}\label{sec:methods-repro}

The pipeline consists of six R modules (data curation, descriptives,
frequentist analysis, subgroups, trial design, simulation) and four
Python modules (machine learning, Bayesian inference, heterogeneity,
RMST), orchestrated by a single \texttt{Makefile} invocation. All scripts
use fixed seeds (global seed 20260916), communicate exclusively through
files in \texttt{data/} and \texttt{outputs/}, and embed
quality-control assertions (record counts, within-patient covariate
consistency, endpoint logic). Figures are produced as vector PDF for the
paper and 300-dpi PNG for the repository README. Exact package versions
are archived with the repository, and every number in this paper was
regenerated and verified from the committed outputs during preparation
of the manuscript.
